% Options for packages loaded elsewhere
\PassOptionsToPackage{unicode}{hyperref}
\PassOptionsToPackage{hyphens}{url}
\PassOptionsToPackage{space}{xeCJK}
%
\documentclass[
  11pt,
  a4paper,
]{ctexart}
\usepackage{amsmath,amssymb}
\usepackage{setspace}
\usepackage{iftex}
\ifPDFTeX
  \usepackage[T1]{fontenc}
  \usepackage[utf8]{inputenc}
  \usepackage{textcomp} % provide euro and other symbols
\else % if luatex or xetex
  \usepackage{unicode-math} % this also loads fontspec
  \defaultfontfeatures{Scale=MatchLowercase}
  \defaultfontfeatures[\rmfamily]{Ligatures=TeX,Scale=1}
\fi
\usepackage{lmodern}
\ifPDFTeX\else
  % xetex/luatex font selection
  \ifXeTeX
    \usepackage{xeCJK}
    \setCJKmainfont[AutoFakeBold=true,AutoFakeSlant=true]{Noto Serif CJK SC}
          \setCJKsansfont[AutoFakeBold=true,AutoFakeSlant=true]{Noto Sans CJK SC}
          \fi
  \ifLuaTeX
    \usepackage[]{luatexja-fontspec}
    \setmainjfont[AutoFakeBold=true,AutoFakeSlant=true]{Noto Serif CJK SC}
  \fi
\fi
% Use upquote if available, for straight quotes in verbatim environments
\IfFileExists{upquote.sty}{\usepackage{upquote}}{}
\IfFileExists{microtype.sty}{% use microtype if available
  \usepackage[]{microtype}
  \UseMicrotypeSet[protrusion]{basicmath} % disable protrusion for tt fonts
}{}
\makeatletter
\@ifundefined{KOMAClassName}{% if non-KOMA class
  \IfFileExists{parskip.sty}{%
    \usepackage{parskip}
  }{% else
    \setlength{\parindent}{0pt}
    \setlength{\parskip}{6pt plus 2pt minus 1pt}}
}{% if KOMA class
  \KOMAoptions{parskip=half}}
\makeatother
\usepackage{xcolor}
\usepackage[margin=2.2cm]{geometry}
\usepackage{longtable,booktabs,array}
\usepackage{calc} % for calculating minipage widths
% Correct order of tables after \paragraph or \subparagraph
\usepackage{etoolbox}
\makeatletter
\patchcmd\longtable{\par}{\if@noskipsec\mbox{}\fi\par}{}{}
\makeatother
% Allow footnotes in longtable head/foot
\IfFileExists{footnotehyper.sty}{\usepackage{footnotehyper}}{\usepackage{footnote}}
\makesavenoteenv{longtable}
\setlength{\emergencystretch}{3em} % prevent overfull lines
\providecommand{\tightlist}{%
  \setlength{\itemsep}{0pt}\setlength{\parskip}{0pt}}
\setcounter{secnumdepth}{-\maxdimen} % remove section numbering
\ifLuaTeX
\usepackage[bidi=basic]{babel}
\else
\usepackage[bidi=default]{babel}
\fi
\babelprovide[main,import]{chinese}
% get rid of language-specific shorthands (see #6817):
\let\LanguageShortHands\languageshorthands
\def\languageshorthands#1{}
\providecommand{\llbracket}{\mathopen{[\![}}
\providecommand{\rrbracket}{\mathclose{]\!]}}
\setlength{\parindent}{2em}
\setlength{\parskip}{0.35em}
\renewcommand{\arraystretch}{1.12}
\ifLuaTeX
  \usepackage{selnolig}  % disable illegal ligatures
\fi
\IfFileExists{bookmark.sty}{\usepackage{bookmark}}{\usepackage{hyperref}}
\IfFileExists{xurl.sty}{\usepackage{xurl}}{} % add URL line breaks if available
\urlstyle{same}
\hypersetup{
  pdftitle={接受之后：基于 eBPF 校准的残差语言与怪异机器论断图谱（中文阅读版）},
  pdfauthor={Chengao Zhang},
  pdflang={zh-CN},
  hidelinks,
  pdfcreator={LaTeX via pandoc}}

\title{接受之后：基于 eBPF 校准的残差语言与怪异机器论断图谱}
\usepackage{etoolbox}
\makeatletter
\providecommand{\subtitle}[1]{% add subtitle to \maketitle
  \apptocmd{\@title}{\par {\large #1 \par}}{}{}
}
\makeatother
\subtitle{独立研究者 · emtanling@gmail.com}
\author{Chengao Zhang}
\date{}

\begin{document}
\maketitle
\setstretch{1.15}

\begin{quote}
\textbf{版本说明（2026-07-17）：}
本文件是供作者核对论证的中文阅读版，不是逐句翻译，也不是投稿规范源。唯一投稿规范源是英文
\texttt{PAPER\_REPORT.tex} 及其编译
PDF；如中文表述与英文规范源不一致，以英文为准。
\end{quote}

\hypertarget{ux6458ux8981}{%
\section{摘要}\label{ux6458ux8981}}

语言理论安全把输入边界建模为形式语言识别器，把下游处理建模为解释。在程序验证器边界，验证器是程序制品的外层识别器；已接受制品又可解释数据并驱动运行时操作词，而任何已计算报告都是另行明确声明的抽象。本文提出五节点论断图谱，严格区分制品接受（A）、同后缀因果状态区分（C）、有界可编程性（P）、由输出见证的已计算报告不可因子分解性（R）以及策略/威胁义务（W）。未来观测商与行为因子分解判据区分以已接受制品为索引的因果语言和相对于报告实例的残差语言。该判据只在选定上下文纤维上的\textbf{唯一单元报告映射}成立时适用；本文不声称每个安全健全但不完备的验证器都必然承载怪异机器。

本文给出统一可控、可观察、可重置门以及固定已接受解释器的证明义务。在精确调度、可接纳性、串行化、源码到目标对应与帧保持前提下，前缀归纳给出有界组合结果。在
Linux eBPF 校准案例中，一个专用、预分配、非 LRU、容量为 2 的哈希映射实现
NAND：重置后保留一个哨兵项，输入比特选择更新已有键或插入新键，第二次更新的成功谓词作为输出。一个固定已接受程序解释输入至多
64 个、门至多 512 个、活跃规范导线至多 578 条的 NAND DAG。

源码快照化的 Linux/aarch64
运行覆盖具名与随机电路、联合边界、串行复用、机制对照和畸形描述符。作者另行运行的语义审计器重建描述符与电路语义，自行签发的清单检查证据包完整性。这个解释器载体记录
A；在声明的具体映射服务与无干扰契约下给出条件性 C
见证；在更多实现前提下支持 P，但不建立 R 或
W。固定辅助可执行实例针对自身计算出的报告建立
\(R(M_{\mathit{linux\_r\_aux\_v1}})\)。另一个已接受 XDP 对象上的本地
stock-Linux/aarch64 实验捕获 exact level 0 的 \texttt{states\_equal}
成功检查及其后的 \texttt{is\_state\_visited}
剪枝，把两条已记录前缀绑定到同一个保留状态代表，并在同一后缀下记录不同观测。把这些记录嵌入事后声明、受哈希绑定的两历史载体和作者定义的操作剪枝报告后，该元组满足定义
2。本文把这一结果表述为``轨迹局部、相对于操作剪枝报告的证书''，而不是
Linux 已规定完整功能报告的证据，也不是覆盖一般 Linux
执行的定理。各载体不可合并：两个 R 结果都没有与解释器的 P
证据链接，也没有任何结果建立 W 或怪异机器。

\textbf{关键词：}
语言理论安全，识别器，残差语言，怪异机器，程序验证，抽象解释，eBPF。

\begin{center}\rule{0.5\linewidth}{0.5pt}\end{center}

\hypertarget{1-ux5f15ux8a00}{%
\section{1. 引言}\label{1-ux5f15ux8a00}}

语言理论安全把输入边界视为语言识别边界：下游处理是解释，输入提供被解释的程序
{[}1{]}--{[}4{]}。本文把这一识别器形状用于程序验证。问题从``接受之后''开始：验证器接受一个程序制品，但该程序还可驱动有状态辅助函数与服务操作。若要把这类下游行为称为残差的、可编程的、由形状诱导的或怪异机器，分别需要什么证据？在这一分层视角下，接受、解释和报告抽象分别回答
\(L_V\) 中的成员关系、\(I\) 中的行为以及 \(\mathsf{Report}_V\)
如何分组这三个不同问题。

制品语言和接受后的操作语言具有不同载体。验证器可以认证有界执行、辅助函数类型使用和内存访问纪律，而不必认证通过每个已文档化运行时服务实现的完整关系。已接受程序可以选择操作、保留状态、观察返回、重置和组合效果；这些事实本身既不证明验证器缺陷，也不证明报告遗漏。

\begingroup
\small
\setlength{\tabcolsep}{2.5pt}
\begin{longtable}[]{@{}p{0.18\textwidth}p{0.42\textwidth}p{0.32\textwidth}@{}}
\toprule\noalign{}
载体 & A/C/P/R/W 状态 & 证据解释 \\
\midrule\noalign{}
\endhead
\bottomrule\noalign{}
\endlastfoot
\texttt{wm\_circuit} 解释器 & A 已记录；C 在声明的服务/无干扰契约下有条件成立；P 在源码/对象、串行化与帧前提下得到支持；R、W 未建立 & 构造性论证与经审计回归证据；解释器前沿没有已计算单元提取器 \\
\(M_{\mathit{linux\_r\_aux\_v1}}\) & A、C、R 在有限自定义元组内成立；P、W 未建立 & 可执行有限模型证书；R 仅针对其计算出的自定义报告，不声称与 stock Linux 存在 refinement 或 bisimulation \\
\(M_K\)（\texttt{rac\_single}） & A 针对冻结对象/内核元组已记录；C 仅在声明的两历史关系中得到见证；R 针对作者定义的操作剪枝报告得到事后证书；P、W 未建立 & 真实剪枝捕获、规范化哈希绑定检查及论文级一步包装；只属轨迹局部，不是 Linux 功能报告契约 \\
\end{longtable}
\endgroup

每一行都有独立的制品、执行载体、报告接口和证据类型。\texttt{wm\_circuit}
行承载本文的 A/C/P 构造；辅助行只针对自定义的报告生成识别器提供正面 R
结果；\(M_K\) 行只针对声明的 \texttt{rac\_single}
操作剪枝报告元组提供事后、轨迹局部的 R
证书。不得跨行组合不同节点，任何已评估行都没有建立 W。

节点 C 和 R 不可共用一个名称。以已接受制品为索引的同后缀语言记为
\(L_{\mathrm{causal}}\)；只有实际计算单元在全部可接纳条件下发生碰撞，词才进入
\(L_{\mathrm{res}}^R\)。P 与 R 在 C 后分叉。策略级怪异机器还要求
P、W、同一编码计算上的链接（Link）与非预期性（Unint）；更窄的``由契约形状诱导、相对于识别器''的分类还要求
R、已文档化且安全保持的语义（Doc）、报告实现符合其声明抽象（Conf）及碰撞确由声明粒度强制（Gran）。该图谱既支持正面证书，也支持有原则的非分类结论。

本文贡献为：（1）带类型的``识别器---解释器---报告''A/C/P/R/W
论断图谱；（2）包含 E1--E4-D 的有界组合证明义务；（3）一个以
\(D_{64,512}\) 为域的 eBPF NAND 解释器校准案例；（4）严格分载体的 R
证据。固定辅助实例针对自定义报告检查有限报告单元碰撞及四个负对照。独立的
stock-Linux 捕获记录真实 exact-0
剪枝以及同一后缀下的不同结果；只有在事后嵌入第 5.7 节显式有限的
\texttt{rac\_single} 载体和作者定义的操作剪枝报告之后，它才构成 R
实例。因此本文只主张该冻结元组的轨迹局部操作证书，而不主张一般 Linux
报告结果；两个 R 实例都不与解释器的 P
证据链接。实证范围限于一个特权、离线 Linux/aarch64
内核构建、两个承载主张的已接受 eBPF
制品、三个额外的已接受解释器控制变体以及一个由自定义识别器接受的辅助制品；不是一般
Linux 报告不透明性定理、并发部署、无界机器、漏洞或策略级怪异机器。

\hypertarget{2-ux76f8ux5bf9ux4e8eux8bc6ux522bux5668ux7684ux6b8bux5deeux8bedux8a00}{%
\section{2.
相对于识别器的残差语言}\label{2-ux76f8ux5bf9ux4e8eux8bc6ux522bux5668ux7684ux6b8bux5deeux8bedux8a00}}

\hypertarget{21-ux8bc6ux522bux6267ux884cux4e0eux62a5ux544a}{%
\subsection{2.1
识别、执行与报告}\label{21-ux8bc6ux522bux6267ux884cux4e0eux62a5ux544a}}

令 \(V\) 为程序制品识别器，\(I\) 为具体执行语义，状态载体为
\(\Sigma_I\)，轨迹载体为 \(\mathcal T_I\)；\(I\)
包含指令引擎、辅助函数、有状态服务及相关环境。因此
\(\mathsf{Tr}_I(P)\subseteq\mathcal T_I\)，且

\[L_V=\{P\mid V(P)=\mathsf{accept}\}.\]

对可选的允许轨迹集合
\(\mathsf{Safe}\subseteq\mathcal T_I\)，安全健全性是

\[\forall P\in L_V.\ \mathsf{Tr}_I(P)\subseteq\mathsf{Safe}.\]

安全健全性是前提，不能由一次成功加载推出；它也不同于抽象变换器的完备性
{[}5{]}。已接受语言、实际计算报告、传递健全性与完备性是不同判断。

若报告在讨论范围内，令

\[\mathsf{Report}_V(P,\ell)\subseteq A_\ell,
\qquad \gamma_\ell:A_\ell\to\mathcal P(\Sigma_I)\]

分别为前沿 \(\ell\) 实际计算的有限单元集及声明的具体化。覆盖要求

\[\mathsf{Reach}_I(P,\ell)\subseteq
\bigcup_{a^\#\in\mathsf{Report}_V(P,\ell)}\gamma_\ell(a^\#).\]

只有同一个实际计算单元同时包含两个状态，二者才是``共同覆盖''；相似日志文本不足以证明单元身份、具体化或共同覆盖。

\hypertarget{22-ux56e0ux679cux8bcdux4e0eux89c2ux6d4bux5951ux7ea6}{%
\subsection{2.2
因果词与观测契约}\label{22-ux56e0ux679cux8bcdux4e0eux89c2ux6d4bux5951ux7ea6}}

对 \(P\in L_V\) 和前沿 \(\ell\)，令 \(\Sigma_{\mathrm{op}}(P)\)
是程序可控运行时操作的有限字母表，\(W^\ell_{\mathrm{run}}(P)\subseteq\Sigma_{\mathrm{op}}(P)^*\)
是代码可从该前沿执行的词。节点 C
使用同后缀测试，不把所有普通运行轨迹重新命名为残差。

比较前固定带类型的观测契约

\[K_{\mathrm{obs}}=(\rho_{\mathrm{obs}},\mathsf{Obs},\mathsf{Slice},\mathsf{Env}),\]

其中 \(\rho_{\mathrm{obs}}:\Sigma_I\to R_{\mathrm{obs}}\)
选择候选状态，\(\mathsf{Obs}:\mathcal T_I\to O_{\mathrm{obs}}\)
投影完整轨迹，\(\mathsf{Slice}\) 为每个词 \(w\) 指定上下文投影
\(\mathsf{ctx}_w:\Sigma_I\to C_w\)，\(\mathsf{Env}\)
是固定环境实例集合。上下文必须包含后缀或观察者读取的全部非选定分量；环境实例固定资源配置、调度、非确定性与外部干扰选择。

写作 \(\llbracket w\rrbracket_{I,e}(\sigma)=(\tau,\sigma')\)。契约在
\(X\subseteq\Sigma_I\) 上对 \(w\) 健全，当任意 \(e\in\mathsf{Env}\) 和
\(\sigma,\sigma'\in X\) 若满足

\[\rho_{\mathrm{obs}}(\sigma)=\rho_{\mathrm{obs}}(\sigma'),\qquad
\mathsf{ctx}_w(\sigma)=\mathsf{ctx}_w(\sigma'),\]

则两次执行有相同的有定义性；若均有定义，则
\(\mathsf{Obs}(\tau)=\mathsf{Obs}(\tau')\)。没有这一非干扰条件，不接纳 C
见证。

\textbf{定义 1（因果状态介导词族）。} 词 \(w\) 在 \((P,\ell)\)
处因果，当
\(P\in L_V\)、\(w\in W^\ell_{\mathrm{run}}(P)\)、\(K_{\mathrm{obs}}\) 在
\(\mathsf{Reach}_I(P,\ell)\) 上对 \(w\) 健全，并存在同一
\(e\in\mathsf{Env}\) 及
\(\sigma_0,\sigma_1\in\mathsf{Reach}_I(P,\ell)\)，使两次后缀执行有定义、终止且

\[
\mathsf{ctx}_w(\sigma_0)=\mathsf{ctx}_w(\sigma_1),\quad
\rho_{\mathrm{obs}}(\sigma_0)\ne\rho_{\mathrm{obs}}(\sigma_1),\quad
\mathsf{Obs}(\tau_0)\ne\mathsf{Obs}(\tau_1).
\]

定义依赖标签族

\[L_{\mathrm{causal}}(V,I;K_{\mathrm{obs}})=
\{(P,\ell,w)\mid w\text{ 在 }(P,\ell)\text{ 处因果}\}.\]

制品与前沿是类型标签。宿主描述符不是 \(L_V\)
的另一个元素，也不会自动成为运行时词；它先成为配置，由固定解释器诱导调度，只有同后缀见证进入
\(L_{\mathrm{causal}}\)。

\hypertarget{23-ux884cux4e3aux5546ux4e0eux62a5ux544aux76f8ux5bf9ux6b8bux5dee}{%
\subsection{2.3
行为商与报告相对残差}\label{23-ux884cux4e3aux5546ux4e0eux62a5ux544aux76f8ux5bf9ux6b8bux5dee}}

固定确定性规约

\[D=(X_D,S_D,A_D,O_D,\delta_D,\lambda_D,s_D,\mathsf{Obs}_D),\]

其中 \(X_D\subseteq\Sigma_I\) 是可接纳具体区域，\(S_D\)
是状态载体，\(A_D\) 是操作字母表，\(O_D\)
是输出字母表，\(\mathsf{Obs}_D:O_D^*\rightharpoonup O_{\mathrm{obs}}\)，并且

\[\delta_D:S_D\times A_D\rightharpoonup S_D,
\qquad\lambda_D:S_D\times A_D\rightharpoonup O_D\]

是同定义域偏函数。对每个使用规约的 \((P,\ell)\)，固定单射操作编码
\(\iota_{P,\ell}:A_D\to\Sigma_{\mathrm{op}}(P)\) 并同态扩展至词。

当且仅当对每个 \(\sigma\in X_D\) 和
\(a\in A_D\)，具体编码步骤存在恰好对应 \(\delta_D(s_D(\sigma),a)\)
有定义，并且每个 \(\sigma\xrightarrow{\iota_{P,\ell}(a)/o}_I\sigma'\)
满足

\[\lambda_D(s_D(\sigma),a)=o,
\qquad s_D(\sigma')=\delta_D(s_D(\sigma),a),
\qquad \sigma'\in X_D,\]

称 \((s_D,\iota_{P,\ell})\) 在 \(X_D\)
上\textbf{操作充分}。因此编码是语义重命名，不能让同一投影状态隐藏不同有定义性、完整输出或后继投影；影响转移的环境选择必须被固定或纳入
\(S_D\)。

令 \(\mathsf{Out}_D(r,w)\) 为迭代 \((\delta_D,\lambda_D)\)
得到的偏完整输出词，\(\mathsf{Def}_D(r,w)\) 表示其有定义，并取
\(\mathcal W_D=A_D^*\)。定义

\[
r\sim_D r'\Longleftrightarrow
\forall w\in\mathcal W_D.\bigl[
\mathsf{Def}_D(r,w)\Leftrightarrow\mathsf{Def}_D(r',w)
\bigr]\land
\bigl[
\mathsf{Def}_D(r,w)\Rightarrow
\mathsf{Out}_D(r,w)=\mathsf{Out}_D(r',w)
\bigr].
\]

\(\sim_D\) 是右同余；若 \(Q_D=S_D/{\sim_D}\) 有限，则 \(X_D\)
上执行诱导商状态上的偏 Mealy 转导器
{[}6{]}--{[}9{]}。本文只测试一个已接受制品、一个前沿上的实际计算单元，不主张提出新的通用完备性理论。

对具体集合 \(F\)，共同启用延续为

\[\mathcal W_D(F)=\{w\in\mathcal W_D\mid
\forall\sigma\in F.\ \mathsf{Out}_D(s_D(\sigma),w)\downarrow\}.\]

\textbf{命题 1（上下文纤维上的行为因子分解）。} 固定
\(P,\ell,D,K_{\mathrm{obs}}\) 和非空
\(F\subseteq\mathsf{Reach}_I(P,\ell)\cap X_D\)。要求：

\begin{enumerate}
\def\labelenumi{\arabic{enumi}.}
\tightlist
\item
  \(\iota_{P,\ell}(\mathcal W_D(F))\subseteq W^\ell_{\mathrm{run}}(P)\)；
\item
  观测契约对全部编码共同延续健全；
\item
  \(F\) 上全部对应上下文投影相同；
\item
  观察者兼容：若编码词 \(w\) 从 \(\sigma\) 执行产生具体轨迹 \(\tau\)，则
  \[\mathsf{Obs}(\tau)=\mathsf{Obs}_D(\mathsf{Out}_D(s_D(\sigma),w));\]
\item
  \textbf{唯一单元条件}：
  \[\forall\sigma\in F.\ \exists!a^\#\in\mathsf{Report}_V(P,\ell).
  \sigma\in\gamma_\ell(a^\#).\]
\end{enumerate}

令 \(\pi_R:F\to\mathsf{Report}_V(P,\ell)\)
把状态映射到其唯一实际计算单元，\(\beta_D(\sigma)=[s_D(\sigma)]_D\)，\(F_{a^\#}=F\cap\gamma_\ell(a^\#)\)。则

\[\exists h:\pi_R(F)\to Q_D.\ \beta_D=h\circ\pi_R\]

当且仅当

\[\forall a^\#\in\pi_R(F).\ |\beta_D(F_{a^\#})|\le1.\]

所以，报告在 \(F\)
上不可因子分解，当且仅当一个唯一实际计算单元包含来自两个不同 \(\sim_D\)
类的状态。唯一单元条件是使 \(\pi_R\)
成为函数的必要前提；若报告标签重叠，必须另行声明成员签名映射，不能沿用本命题。

记 \(\mathsf{Adm}(P,\ell,D,F;K_{\mathrm{obs}})\)
表示接受、操作充分、可达非空纤维、运行词包含、观测者兼容、健全观测契约、共同上下文和唯一单元条件全部成立。唯一单元条件使非空的
\(F_{a^\#}\)
构成不交覆盖；单元重叠的报告需要另行声明成员签名映射，不属于本判据。

可接纳性是内部语义定型条件，并不自动排除事后选择。每个正面实例必须声明
\(D\)、\(F\)、观测者、后缀和报告接口究竟是在看到区分性执行之前预先确定，还是在结果已知后事后选定。若要提出前瞻性或报告层面的一般主张，还须提供独立于见证输出的选择来源，例如预先登记的协议或外部规定的报告契约。事后元组可以严格满足定义
2，但仅凭这一事实不能刻画识别器意图中的报告、更广的执行域或一类实现。

\textbf{定义 2（相对于报告的输出见证残差语言）。}
对一个可接纳元组，带标签词 \((P,\ell,D,F,a^\#,w)\)
属于输出见证的报告相对残差，当：编码词属于
\(L_{\mathrm{causal}}\)；\(w\in\mathcal W_D(F)\)；定义 1 的两个见证位于
\(F\)；二者唯一报告标签同为 \(a^\#\)；且其 \(\beta_D\)
类不同。\(L_{\mathrm{res}}^R\) 是所有可接纳元组上这类带标签词的并集。

该定义只是报告因子分解失败的\textbf{输出见证子集}。商关系也区分有定义性，因此可能存在不可因子分解，却没有一个两边都终止并产生不同输出的词。保留的
\texttt{wm\_circuit} Linux
日志不提供实际计算单元提取器、具体化或唯一单元映射；所以该解释器载体只给出条件性
\(L_{\mathrm{causal}}\) 见证，不建立 \(L_{\mathrm{res}}^R\)。第 5.7
节为另一个有限的 \texttt{rac\_single} 载体实例化这些对象。

\hypertarget{24-ux5f62ux72b6ux7f3aux9677ux4e0eux7b56ux7565}{%
\subsection{2.4
形状、缺陷与策略}\label{24-ux5f62ux72b6ux7f3aux9677ux4e0eux7b56ux7565}}

缺陷诱导缝隙违反声明的识别器、报告或运行时契约。一个使用已文档化且安全保持语义、并在报告意图认证的关系上具有
\(L_{\mathrm{res}}^R\) 见证的案例，只是契约形状缝隙的证据；仍须以 Conf
与 Gran
排除错误变换器、汇合或提取器。策略级怪异机器还需同一计算上的行为主体控制、策略排除效果与非预期解释。普通有状态
API 即使实现转导器，也未必满足任一分类。

\hypertarget{3-ux6709ux754cux72b6ux6001ux4ecbux5bfcux7684ux7535ux8defux5b9eux73b0}{%
\section{3.
有界状态介导的电路实现}\label{3-ux6709ux754cux72b6ux6001ux4ecbux5bfcux7684ux7535ux8defux5b9eux73b0}}

\hypertarget{31-e1e3ux4eceux533aux5206ux5230ux53efux590dux7528ux95e8}{%
\subsection{3.1
E1--E3：从区分到可复用门}\label{31-e1e3ux4eceux533aux5206ux5230ux53efux590dux7528ux95e8}}

固定规约 \(D\)、规范重置商类 \(q_0\in Q_D\) 及非空
\(\mathsf{Adm}_G\subseteq X_D\)。门基
\((P_G,\mathit{reset},G,\mathit{observe},D,\mathsf{Adm}_G)\)
的操作编码在 \(X_D\) 上操作充分，并满足：

\begin{itemize}
\tightlist
\item
  \textbf{E1（因果基与观测）}：存在输入 \(x_0,x_1\)、前缀
  \(p_0,p_1\)、同一剩余后缀 \(w\) 和内部前沿 \(\ell_G\)，使
  \(u_G(x_i)=p_iw\)；两个前缀所达状态是定义 1 对编码 \(w\)
  的见证；完整门读出恰等于该因果观测，并正是 E4-D 写入的门结果。
\item
  \textbf{E2（统一输入控制）}：同一 \(P_G\) 中的分派器读取
  \(x\in\{0,1\}^2\)。从 \(q_0\) 中任意可接纳状态执行 \(u_G(x)\)
  都有定义，且
  \[\mathit{observe}(\mathsf{Out}_D(s_D(\sigma),u_G(x)))=g(x).\]
\item
  \textbf{E3（重置）}：重置词对所有可接纳状态有定义，保持导线单元并返回
  \(q_0\)；固定上下文和导线相同的两个重置结果，对所有输入产生相同完整输出词。
\end{itemize}

E1 防止读出与状态区分无关；E2 防止实验者从外部提供真值表；E3
防止通道只能使用一次。

\hypertarget{32-ux63cfux8ff0ux7b26ux57dfux4e0e-e4-d}{%
\subsection{3.2 描述符域与
E4-D}\label{32-ux63cfux8ff0ux7b26ux57dfux4e0e-e4-d}}

域 \(D_{64,512}\) 的描述符为

\[d=(m,n,(s_i^0,s_i^1)_{0\le i<n}),\]

其中 \(0\le m\le64\)、\(0\le n\le512\) 且 \(0\le s_i^b<2+m+i\)。导线
0、1 为常量，输入位于 2 至 \(2+m-1\)，门 \(i\) 写规范目的
\(2+m+i\)；最多 578 条活跃规范导线，最高索引 577。

\(\mathsf{Eval}_g(d,x)\) 按描述符顺序扩展完整规范导线向量：

\[\nu[2+m+i]=g(\nu[s_i^0],\nu[s_i^1]).\]

\(\mathsf{Enc}_U(d,x)\) 把规范描述符、输入和控制写入映射；WMC1
只是宿主序列化，不是 \(V\) 接受的 BPF 程序。固定解释器 \(P_U\)
读取配置并诱导 \(\mathsf{Sched}_U(d,x)\)。状态屏蔽接口只在状态为
\(\mathsf{OK}\) 时返回规范导线观察；这不声称物理陈旧单元被擦除。

\textbf{E4-D（有界数据参数化解释）：}要求同一个 \(P_U\) 与 E1--E3
门基共享制品，并满足：固定映射定义与记录环境中的接受独立于
\(d,x\)；有效配置初始化常量和输入，恰按 \(0,\ldots,n-1\)
执行并成功终止；一个外部临界区覆盖全部共享映射的设置、调用和回读；每次迭代验证规范形式、复制两个较早源、重置并调用门，只写规范目的与审计单元并保持描述符和既有导线；畸形配置在至多
512 次内以非 \(\mathsf{OK}\) 终止并返回语义结果 \((s,\bot)\)。

\hypertarget{33-ux5b9eux73b0ux5b9aux7406}{%
\subsection{3.3 实现定理}\label{33-ux5b9eux73b0ux5b9aux7406}}

\textbf{定理 1（显式义务下的有界组合）。} 若同一固定 \(P_U\)
在声明的确定、串行化环境下履行 E4-D，且其门基以函数 \(g\) 履行
E1--E3，则对所有 \(d\in D_{64,512}\) 和 \(x\in\{0,1\}^{m(d)}\)，

\[\mathsf{Run}_{U,d}(\mathsf{Enc}_U(d,x))
=(\mathsf{OK},\mathsf{Eval}_g(d,x))\]

并恰执行 \(n(d)\) 次迭代。若识别边界另有 \(\mathsf{Safe}\)
的安全健全性前提，这些轨迹才可额外推出属于 \(\mathsf{Safe}\)。

证明由规范前缀归纳给出：E4-D
初始化前缀与精确调度，描述符界保证源均已建立，E3 回到 \(q_0\)，E2 给出
\(g\)，E4-D 只写新目的并保持前缀；E1
只负责把该功能位绑定到同后缀因果区分。定理分解证明义务，不声称测试枚举了描述符域，也不证明``每个电路生成一个新
BPF 对象''的 E4-A。

\hypertarget{4-ebpf-ux6821ux51c6ux6848ux4f8b}{%
\section{4. eBPF 校准案例}\label{4-ebpf-ux6821ux51c6ux6848ux4f8b}}

\hypertarget{41-ux6267ux884cux4e0eux8f7dux4f53ux8fb9ux754c}{%
\subsection{4.1
执行与载体边界}\label{41-ux6267ux884cux4e0eux8f7dux4f53ux8fb9ux754c}}

固定程序 \(P_U=\mathit{wm\_circuit}\) 的节为
\texttt{SEC("syscall")}。记录环境中，内核验证器接受该对象，用户态通过
\texttt{bpf\_prog\_test\_run\_opts()} 离线执行；该解释器运行没有实时钩子
{[}10{]}, {[}11{]}。验证单元是固定 \(P_U\)，不是 WMC1
或描述符；接受后，\(P_U\) 读取映射并生成辅助函数操作与导线观察。

门映射 G0 是 \texttt{BPF\_MAP\_TYPE\_HASH}，\texttt{max\_entries=2}，非
LRU，保持默认预分配
{[}12{]}。规约要求重置成功，并由外部临界区覆盖设置、调用与回读；程序本身没有并发锁。

\hypertarget{42-ux9971ux548cux79e9-nand-ux4e0eux6761ux4ef6ux6027-c}{%
\subsection{4.2 饱和秩 NAND 与条件性
C}\label{42-ux9971ux548cux79e9-nand-ux4e0eux6761ux4ef6ux6027-c}}

使用互异的哨兵键 \(S\) 与输入键 \(A,B\)。重置删除三键后插入
\(S\)。零更新 \(S\)，第一个一插入 \(A\)，第二个一插入
\(B\)。在有效参数、默认预分配、成功重置、无干扰和记录 Linux 6.17
环境前提下，已有键更新成功且不增占用，新键在容量以下成功并增占用，容量已满时失败；只观察第二次更新是否成功，不依赖具体错误码
{[}12{]}。四种输出为 \(1,1,1,0\)，即 NAND。NAND 具功能完备性，因为
\(\neg x=\mathsf{NAND}(x,x)\) 且 \(x\land y=\neg\mathsf{NAND}(x,y)\)。

定义 1 的具体见证取内部前沿为第一次输入条件操作之后、第二次更新之前。令
\(R_{G0}(\sigma)\) 为专用映射 G0
的\textbf{完整辅助函数相关动态状态}，包括键占用、桶、预分配元素/空闲链元数据以及更新辅助函数读取的任何其他映射局部状态，并取
\(\rho_{\mathrm{obs}}(\sigma)=R_{G0}(\sigma)\)。占用键集合 \(K_{G0}\)
只是 \(R_{G0}\) 的派生投影，不能单独充当完整选定状态。

输入 \((0,1)\) 与 \((1,1)\) 在该前沿的派生占用分别为 \(\{S\}\) 与
\(\{S,A\}\)，所以完整状态不同。公共后缀 \(w\) 是插入新键
\(B\)，随后统一捕获并存储成功位；\(\mathsf{Obs}(\tau)=[\mathit{ret}(\tau)=0]\)。\(\mathsf{ctx}_w\)
包含程序点、键和值、G0
身份与静态属性、模式、相关控制/导线值，以及所有位于 \(R_{G0}\)
之外但被后缀读取的分量；\(\mathsf{Env}\)
固定内核、对象、映射实例、调度及无干扰。在声明的映射更新契约下，两次观察为
1 与 0。较早输入不被后缀读取，其持续且与后缀相关的效应包含于完整
\(R_{G0}\) 中。因此这是\textbf{依赖具体服务契约的条件性 C
见证}，不是内核内部状态已被轨迹完整暴露的主张。

门规约的 \(s_{D_G}(\sigma)=(\mathit{phase}(\sigma),K(\sigma))\)
只在由有效参数、专用预分配映射契约、成功重置、串行化与无干扰界定的
\(X_{D_G}\)
上使用。其操作充分性是显式服务契约前提，不是对所有内核字段的机器检查模型。E1
由上述内部前沿见证支持；命题 2 与重置前提支持 E2--E3。该机制不增加 eBPF
的外延布尔表达能力，也不单独证明报告遗漏。

\hypertarget{43-ux56faux5b9aux89e3ux91caux5668ux524dux7f00ux4e0dux53d8ux5f0fux4e0eux8bc1ux636eux8fb9ux754c}{%
\subsection{4.3
固定解释器、前缀不变式与证据边界}\label{43-ux56faux5b9aux89e3ux91caux5668ux524dux7f00ux4e0dux53d8ux5f0fux4e0eux8bc1ux636eux8fb9ux754c}}

解释器还使用
TAPE、CIRCUIT、WIRES、VM\_CONTROL、VM\_TRACE。宿主解析并规范化
WMC1，写映射，只在 \texttt{status=OK} 后投影请求输出；输出列表不由 BPF
读取。内核中有界循环处理至多 512
个描述符；每次迭代先复制描述符和源值，再调用辅助函数、重置
G0、执行门并按键更新规范目的，避免跨辅助函数保留映射值指针。

源码级前缀不变式在下列前提下成立：命题 2
更新律、成功重置及必需调用、E4-D
串行化、有效编码，以及捕获目标保留人工检查的源码控制/数据依赖。\texttt{bpf\_loop}
请求 \(n(d)\) 次迭代；回调返回 0 继续、1 停止，辅助函数返回执行次数
{[}13{]}；解释器只有在该返回值与自身完成计数均为 \(n(d)\)
时接受。归纳得到每个规范前缀等于
\(\mathsf{Eval}_{\mathsf{NAND}}\)。这是源码级论证加人工检查转储，不是机器检查的
eBPF 语义证明。

因此：A 已记录；C 仅在声明的映射服务/无干扰契约下为条件性；P
还依赖源码到对象、帧保持和串行化前提。对这个 \texttt{wm\_circuit}
载体，没有解释器前沿的 Linux 报告单元提取器或部署策略，故其 R 与 W
均未建立。

\hypertarget{5-ux8bc4ux4f30}{%
\section{5. 评估}\label{5-ux8bc4ux4f30}}

主要证据位于
\texttt{results/interpreter/interpreter-final-20260711-02/}，环境为
Ubuntu 24.04、Linux 6.17.0-35-generic、aarch64。证据包保留四个 BPF
变体、已加载程序元数据、验证器日志、描述符、JSONL、作者运行的语义审计、当时选择的源码/手稿快照及自行签发的
SHA-256
清单。最终手稿晚于该运行，不受该清单覆盖；证据由作者生成并由作者另行审计，不是第三方复现。

证据行总数为

\[38{,}533=26{,}488\text{ 条逐门记录}+12{,}037\text{ 条成功运行记录}+8\text{ 条负对照记录},\]

它们是异构证据行，不是``38,533 个测试''。覆盖包括：9 个具名电路（39
次成功、166 门记录）；100 个固定种子随机 DAG（1,876 次、23,776
门记录）；512 门深度与 64 输入/512 门/578 线联合边界（4 次、2,048
门记录）；零门与 10,000 次串行调用（10,001 次）；三种机制/基线对照（117
次、498 门记录）；8 个畸形案例。联合 64 输入边界只测试全零和全一，不是
\(2^{64}\) 穷举；整个语料库是回归证据，不是 \(D_{64,512}\) 的枚举证明。

独立语义审计重建具名和边界描述符、再生成随机语料、独立计算完整导线向量、检查目的与畸形案例并匹配运行标签；审计报告成功。清单验证也成功，但清单不是签名、时间戳、认证或独立复现。

容量 64 与强制哨兵对照均使具名语料 166 个门输出全为
1；算术基线产生预期结果。它们支持``容量饱和和第二个新键共同导致零输出''的机制归因，不证明等价验证器报告。Linux
验证器文档的目标是程序安全、路径、参数和内存访问
{[}14{]}，不承诺映射介导计算的完备功能证书。保留的 \texttt{wm\_circuit}
证据在解释器前沿没有报告提取器，因此不推断验证器单元跟踪该程序辅助函数返回值的精度。该载体用于校准
LangSec 中``已识别安全性''与``被解释功能''之间的边界：其 R
仍未建立，但这不影响第 5.7
节以专用已接受对象和显式操作剪枝报告契约评估另一个 stock-Linux 载体。

\hypertarget{56-ux56faux5b9aux8f85ux52a9ux53efux6267ux884cux62a5ux544aux5b9eux4f8b}{%
\subsection{5.6
固定辅助可执行报告实例}\label{56-ux56faux5b9aux8f85ux52a9ux53efux6267ux884cux62a5ux544aux5b9eux4f8b}}

为在不重命名 Linux 验证器制品的前提下执行报告相对判据，将辅助 R
实例嵌入第 7.1 节的完整模型载体并固定

\[M_{\mathit{linux\_r\_aux\_v1}}=
(V_{\mathit{linux\_r}},I_{\mathit{hash}},\mathsf{Report}_{\mathit{aux}},
K_{\mathrm{obs}},P_{\mathit{aux}},\ell,D,F,
\varnothing,\varnothing,\varnothing,\varnothing).\]

\(P_{\mathit{aux}}\) 由自定义报告生成识别器 \(V_{\mathit{linux\_r}}\)
接受；它不是经 stock Linux 验证器接受的 BPF 对象。\(I_{\mathit{hash}}\)
是有限、确定、串行且无干扰的受限语义，只覆盖固定程序使用的非驱逐 HASH
映射更新情形。这里 \(\mathsf{Report}_{\mathit{aux}}\) 只指
\texttt{report.json{[}"report\_cells"{]}}；\texttt{derivation.json}
仅记录工作列表/计算溯源，不属于报告标签接口。最后四项分别是有类型的空行为主体集、空效果集、空驱动关系和空许可效果集；它们只把
R 实例嵌入共同载体，不提出 W 论断。

\textbf{可执行证书。} \(R(M_{\mathit{linux\_r\_aux\_v1}})\)
成立（制品状态：established）。

\emph{证书论证。} 识别器接受
\(P_{\mathit{aux}}\)。在单一串行无干扰环境中穷尽
\(a\in\{0,1\}\)、\(b=1\)，得到
\(F=\mathsf{Reach}_{I_{\mathit{hash}}}(P_{\mathit{aux}},\ell)
=\{\texttt{frontier:S},\texttt{frontier:AS}\}\subseteq X_D\)；\(X_D\)
还包含相应的两个终止状态。令 \(a_{\mathrm{obs}}\in A_D\) 表示名为
\texttt{update-suffix-and-observe} 的规约动作，以 \(c_{\mathrm{obs}}\)
表示精确编码的运行时操作符号
\texttt{bpf\_map\_update\_elem(G0,suffix\_key,one,BPF\_ANY);observe(ret==0)}，并令
\(\iota_{P_{\mathit{aux}},\ell}(a_{\mathrm{obs}})=c_{\mathrm{obs}}\)；这个单元素编码是单射。\(\iota_{P_{\mathit{aux}},\ell}(a_{\mathrm{obs}})\)
的具体执行在两个前沿状态上有定义，在终止状态上无定义，并在定义性、输出和后继上与
\(D\) 一致，故检查器在整个 \(X_D\) 上建立操作充分性。观测契约把
\(\rho_{\mathrm{obs}}\) 固定为已占用键集合，把 \(\mathsf{Obs}\)
固定为有序成功位词，把 \(\mathsf{Slice}\)
固定为程序阶段/服务上下文对，把 \(\mathsf{Env}\)
固定为上述单一环境；特别地，\(\rho_{\mathrm{obs}}(\texttt{frontier:S})=\{S\}\ne\{S,A\}=\rho_{\mathrm{obs}}(\texttt{frontier:AS})\)。对于
\(F^2\) 中每个有序状态对和
\(\mathcal W_D(F)=\{\varepsilon,a_{\mathrm{obs}}\}\)
中每个词，检查器均建立其余运行时词、共同上下文、观察者兼容性与
\(K_{\mathrm{obs}}\)
健全性义务。这些义务履行了除唯一单元覆盖之外的全部可接纳条件。

\(\mathsf{Report}_{\mathit{aux}}\) 发出一个唯一单元
\(a^\#\)，共同覆盖两个前沿状态，从而完成
\(\mathsf{Adm}(P_{\mathit{aux}},\ell,D,F;K_{\mathrm{obs}})\)。精确有限未来观测商把二者分到不同类，故
\(|\beta_D(F_{a^\#})|=2\)，且 \(a_{\mathrm{obs}}\) 分别输出 1 与
0。同一后缀因此先见证定义 1，而带标签元组
\((P_{\mathit{aux}},\ell,D,F,a^\#,a_{\mathrm{obs}})\) 满足定义 2，所以
\(R(M_{\mathit{linux\_r\_aux\_v1}})\) 成立。证据另含 21
项域/动作的\textbf{返回类别包含检查}和 2
项报告单元的\textbf{后继包含检查}，均为零违反；前 21 项不是 21
个完整后状态检查。精确占用跟踪、容量
64、强制哨兵、忽略返回四个负对照均消除 R。\(\square\)

在归档证据包上，另行实现且由作者运行的检查器不导入模型实现，重构可达性、唯一单元覆盖、商和因子分解判定，并给出辅助
R
已建立的判定。另一个回归测试仅把同一确定性模型证据包构建两次并逐字节比较五个形式化
JSON
文件；它是确定性测试，不是实验性证据。该证书属于可执行有限模型证据，不是机器检查证明。保留的内核校准对
\((0,1)\) 与 \((1,1)\) 两个赋值共有 4 行
oracle；它只校准这两个受限服务结果，不证明 \(I_{\mathit{hash}}\) 与
Linux 之间的 refinement 或 bisimulation，也不提取 stock
验证器单元。因此辅助证书与下节 stock-Linux 证书彼此独立。其证据路径为
\texttt{results/linux\_r/linux-r-v1/}。

\hypertarget{57-ux51bbux7ed3-stock-linux-ux64cdux4f5cux526aux679dux62a5ux544aux5b9eux4f8b}{%
\subsection{5.7 事后 stock-Linux
轨迹证书}\label{57-ux51bbux7ed3-stock-linux-ux64cdux4f5cux526aux679dux62a5ux544aux5b9eux4f8b}}

本节把已经完成的本地 stock-Linux
实验作为事后轨迹证书分析。内核捕获是一手证据，证明冻结元组上发生了一个剪枝事件及两个运行结果。集成检查器与冻结包检查器离线验证声明的证据关系；其回归测试不能替代该捕获，也不能把作者定义的操作剪枝报告转化为
Linux 文档规定的功能报告契约。

为在不复用辅助语义的前提下执行定义
2，固定如下受哈希绑定、仅限证据范围的载体：

\[M_K=(V_K,I_K,\mathsf{Report}^{\mathit{prune}}_K,K^K_{\mathrm{obs}},
P_R,\ell_K,D_K,F_K,\varnothing,\varnothing,\varnothing,\varnothing).\]

\(P_R\) 是固定且经验证器接受的 XDP 对象 \texttt{rac\_single}，不是解释器
\(P_U=\mathit{wm\_circuit}\)。\(V_K\) 是记录中的 stock Linux
6.17.0-35-generic 验证器，精确比较级别为
0；\(\ell_K=\operatorname{pre}(41)\) 是共同 \texttt{shared\_suffix}
调用前的调用者侧翻译后前沿。对这个事后命题，\(I_K\)
是在捕获完成后，由论文把证据包中的两次已记录执行限制并加阶段标签而构造的有限串行关系。它既不是证据包适配器的自环转移系统，也不是所有
Linux
执行的集合。其声明的输入/环境域只含两条已捕获选择器历史，所以下面的可达性等式是受限分析载体的定义，而不是关于
Linux 全部可达状态的经验完备性主张：

\[\mathsf{Reach}_{I_K}(P_R,\ell_K)=F_K=\{\sigma_0,\sigma_1\}.\]

令
\(X_{D_K}=\{\sigma_0,\sigma_1,\sigma_0^+,\sigma_1^+\}\)；这些符号表示针对两个前缀案例及其后缀后记录构造的阶段化具体化见证，并非完整运行时寄存器、栈和辅助函数内部映射状态的直接快照。两个
\texttt{runtime.json} 最终状态记录只绑定派生的 G0 键集合投影
\(\{S,B\}\)、\(\{S,A\}\) 及程序级结果；验证器捕获另行绑定选定的 State V2
记录。令 \(a_B\) 为抽象插入 B 动作，并定义宏符号
\(c_B\in\Sigma_{\mathrm{op}}(P_R)\)，其声明的具体解释执行完全相同的剩余翻译后后缀；置
\(\iota_{P_R,\ell_K}(a_B)=c_B\)。一步规约为

\[
\begin{gathered}
S_{D_K}=\{p_0,p_1,q^{\mathrm{term}}_0,q^{\mathrm{term}}_1\},\\
A_{D_K}=\{a_B\},\qquad O_{D_K}=\{0,1\},\\
O^K_{\mathrm{obs}}=\{\varepsilon,0,1\},\\
s_{D_K}(\sigma_i)=p_i,\quad s_{D_K}(\sigma_i^+)=q^{\mathrm{term}}_i,\quad
\delta_{D_K}(p_i,a_B)=q^{\mathrm{term}}_i,\\
\lambda_{D_K}(p_0,a_B)=1,\qquad \lambda_{D_K}(p_1,a_B)=0.
\end{gathered}
\]

在具体层面，声明的关系对该动作恰好包含

\[
\sigma_i\mathrel{\overset{c_B/b_i}{\longrightarrow}_{I_K}}\sigma_i^+,
\qquad (b_0,b_1)=(1,0),
\]

且从 \(\sigma_0^+,\sigma_1^+\) 均无 \(c_B\)
转移。这些转移属于论文定义的包装；证据包中的 \texttt{runtime.json}
记录只绑定键集合投影与程序级结果，不绑定完整具体端点。从
\(q^{\mathrm{term}}_0,q^{\mathrm{term}}_1\) 没有出边。两个观测者均返回
\(O^K_{\mathrm{obs}}\)
中的位词：\(\mathsf{Obs}_{D_K}(\varepsilon)=\varepsilon\) 且
\(\mathsf{Obs}_{D_K}(b_i)=b_i\)；在相应的声明具体侧，\(\mathsf{Obs}\)
把空迹映为 \(\varepsilon\)，把每条已记录 \(c_B\) 迹映为
\(b_i\)。因此观测兼容性覆盖
\(\mathcal W_{D_K}(F_K)=\{\varepsilon,a_B\}\)
中两个词。阶段标签和所列转移使操作充分及运行词包含在这个受限关系中按构造成立；它们并未独立验证更大的
Linux 转移语义。

stock 专用契约 \(K^K_{\mathrm{obs}}\) 置
\(\rho^K_{\mathrm{obs}}(\sigma)=R_{G0}(\sigma)\)，即第 4.2
节定义的完整辅助函数相关 G0
动态状态，并观察程序级成功位词。记录的键集合只是派生投影
\(K_{G0}\)：\(K_{G0}(\sigma_0)=\{S\}\) 且
\(K_{G0}(\sigma_1)=\{S,A\}\)，所以即使记录没有暴露每个内部字段，也有
\(R_{G0}(\sigma_0)\ne R_{G0}(\sigma_1)\)。上下文固定调用方阶段与后缀、内核/对象/程序及映射身份、键值常量、映射类型/容量/标志、串行化、单制品条件以及
\(R_{G0}\)
之外所有被后缀读取的分量。检查器比较声明的七个规范化运行时上下文字段；另由论文审查翻译后
PC 41-\/-44 的调用方调用/返回路径及 PC 107-\/-122 的
\texttt{shared\_suffix}
被调用函数读集。被调用函数会覆盖键/值工作寄存器和栈槽；在辅助函数之外，它除此之外只读取固定常量与固定的
G0
映射身份，而调用方仅调用该函数并返回其结果。映射局部桶、元素/空闲链元数据及辅助函数读取的其他
G0 状态属于 \(R_{G0}\)；哈希绑定的单元素 \(\mathsf{Env}\)
固定内核、对象、精确映射实例、服务/分配选择、调度与无干扰。因此两条记录具有相同的非选定后缀读取上下文。由于选定状态的投影已不同，契约在
\(F_K\) 上``\(R_{G0}\)
相等''的前提只会把一个状态与自身配对；声明的一步服务关系具有确定性，故契约健全。记录的成功位等于
\(\mathsf{Obs}_{D_K}\)，故观察者兼容。这样，运行词、共同上下文、操作充分、观察者兼容和健全性义务只在这个显式有限载体上得到履行，而非在一般
Linux 执行上得到履行。该读集论证以及 \(I_K,D_K,K^K_{\mathrm{obs}}\)
的定型是论文审查的义务，不是对检查器机器验证范围的扩大表述。

\(\mathsf{Report}^{\mathit{prune}}_K\)
是作者为分析该捕获事件而显式选定的\textbf{操作剪枝报告}：\texttt{states\_equal}
成功检查实际导致 \texttt{is\_state\_visited} 剪除 current state
时，它把可接纳前缀见证分配给 retained verifier-state
representative。有向剪枝边就是成员关系；本文不把 \texttt{states\_equal}
重释为完整具体状态上的对称数学等价。Linux
把验证器文档化为安全机制，而没有把这一接口规定为覆盖运行时映射占用的完备功能报告。因此这个报告是从具体剪枝事件提取的分析投影，不是已被实验推翻的
Linux 指定证书。

冻结清单绑定内核版本、BTF 与配置哈希、对象 SHA-256、已加载程序
id/tag/pin 以及翻译后字节码 SHA-256。选定 fexit 事件在指令 41 同时记录
\path{states_equal_success=true} 与
\path{is_state_visited_prune=true}。retained 与 current
的历史不同，分别经过 \(a=1\) 与 \(a=0\)
对应的翻译后分支调用，再到达同一前沿；二者状态哈希也不同（\path{516c47f044cc3fc3}
与
\path{a66d912d58c91de4}），成员关系来自观测到的有向剪枝边，而非哈希相等。经审查的路径相关检查器把两条历史绑定到该前沿和同一剩余字节码后缀；捕获完成记录显示丢失事件与解析错误均为
0。

规范化成员检查把构造见证 \(\sigma_0\) 与捕获的 current verifier state
关联，把 \(\sigma_1\) 与捕获的 retained verifier state
关联；这些检查并未捕获完整运行时前沿状态。运行记录把见证的派生 G0
键集合投影分别绑定为 \(\{S\}\) 与 \(\{S,A\}\)。记
\(a_K^\#=\mathtt{retained{:}516c47f044cc3fc3}\)。在受限 exact-0 State V2
形状和声明的操作报告内，形状检查器与唯一单元检查器建立

\[\gamma^K_{\ell_K}(a_K^\#)\cap F_K
=\{\sigma_0,\sigma_1\}.\]

而且在 \(F_K\) 上，两状态都不属于任何其他已提取报告单元。因此

\[\pi_R(\sigma_0)=\pi_R(\sigma_1)=a_K^\#.\]

在共同定义的词 \(a_B\) 下，记录的\textbf{程序级成功位}从 \(\sigma_0\) 与
\(\sigma_1\) 出发分别为 1 与 0；这两个位由 \texttt{observe(ret==0)}
派生，不是映射更新辅助函数的原始返回值（成功为 0、失败为负值）。根据 \(\sim_{D_K}\)
的定义，输出不同直接蕴含 \(p_0\not\sim_{D_K}p_1\)，故
\(\beta_{D_K}(\sigma_0)=Q_0\ne Q_1=\beta_{D_K}(\sigma_1)\)，而操作报告把两个具体状态分配给同一单元。接受与身份、可达路径、直接成员、会话完整性、唯一单元、商、因子分解和哈希门均通过；结合上面显式定型的一步构造，这些事实建立
\(\mathsf{Adm}(P_R,\ell_K,D_K,F_K;K^K_{\mathrm{obs}})\) 并满足定义
2。因此

\[R(M_K)\]

只对该声明的有限冻结载体成立。本文把它报告为\textbf{事后、轨迹局部、相对于操作剪枝报告的证书}。集成检查器输出
\texttt{STOCK\_LINUX\_R\_ESTABLISHED\_FOR\_FROZEN\_TUPLE}，但该标识只表示为此元组编码的哈希绑定证据门均已通过；它不表示
Linux 暴露或违反了文档规定的完备功能报告。\(I_K,D_K,K^K_{\mathrm{obs}}\)
的定型和上述健全性论证是把一步记录置于定义 2
下的论文级形式包装，不是覆盖完整 Linux
执行空间的机器检查定理。证据包适配器在重算稳定商时把每个前缀案例表示为自环；第一次
\(a_B\) 输出已不同，故细化在深度 1
即分开两状态。用上述终止转移替换自环仍保留该划分及独立提取的操作报告映射，但本文不把两个商模型视为相同。所需
\(\beta\) 不等式与输出见证子句由 \(a_B\) 证明，不导入任何长词结果。

该证书同时相对于载体、元组、观测者、后缀和报告。由于这些选择在区分性捕获出现后才最终确定，它支持声明的事后元组上的存在性结论，但不支持独立预承诺或总体层面的
Linux 主张。它不把 \(I_K\)
推广到未捕获的输入、续延或可达状态，不推广到其他内核、配置、BTF、编译输出或
eBPF 对象，不刻画一般 \texttt{states\_equal}
语义，也不证明验证器不健全、漏洞、P、W
或怪异机器；它也不证明映射占用属于 Linux
指定的完备功能报告。具体化论证只涉及与捕获的 State V2
记录及局部未来可观测字段兼容的构造见证。由于
\(P_R\ne P_U\)，该证书不能与解释器的 P 证书组合以满足
\(\mathsf{Link}\)。冻结证据位于
\texttt{residuality-auditor/stock-linux-r-proof/}。

\hypertarget{58-ux6709ux6548ux6027ux5a01ux80c1}{%
\subsection{5.8 有效性威胁}\label{58-ux6709ux6548ux6027ux5a01ux80c1}}

实验只使用一个内核构建版本和一种体系结构。论证依赖专用的预分配非 LRU
映射、成功重置、互异的键、所述更新律，以及整个映射集合上的互斥。会发出门记录的数据集保留原始返回值；10,000
次运行的压力测试文件记录运行/状态/输出证据，但不包含逐门原始行。证明只使用成功与失败之分，并不承诺可移植的错误码。

有限测试不能证明所有 \(D_{64,512}\)。定理 1
转而依赖所陈述的源代码级初始化、验证、有序迭代、门和帧义务；语料库旨在寻找反例。翻译后转储和验证器日志被保留下来以供人工检查及清单哈希使用，而未被自动化数据流检查器消费。状态掩蔽具有源代码顺序保护，而负面测试套件测试的是拒绝/状态，并非注入的掩蔽路径。不存在机器检查的
eBPF 语义证明，且只有在可选的 \(\mathsf{Safe}\)
结论中才假定生产验证器的安全健全性。

stock-Linux
实验是在一个冻结内核/对象元组上由作者生成并审查，尚未得到第三方独立复现。操作报告、两历史载体、观测者和后缀都是事后最终确定，因此即使内部检查全部通过，结果仍受选择效应威胁。证据包绑定
BTF、配置与源码映射元数据，但没有保存目标 \texttt{verifier.c}
函数正文；受限形状引理不是 Linux 一般 \texttt{states\_equal}
语义的源码级定理。

原始捕获阶段分析器有意保持
fail-closed：在字节码前沿、路径对应、报告契约和具体化关系尚未审查时，它仍输出
\texttt{LINUX\_R\_NOT\_ESTABLISHED\_FROM\_RAW\_CAPTURE}。受哈希绑定的集成检查器随后消费这些经审查证书并给出元组特定的证据判定；正式
R 结论还依赖第 5.7
节有限定型包装。仅凭原始事件、状态哈希或类似验证器日志文本不能满足定义
2。

\hypertarget{6-ux76f8ux5173ux5de5ux4f5c}{%
\section{6. 相关工作}\label{6-ux76f8ux5173ux5de5ux4f5c}}

语言理论安全提供识别/解释框架
{[}1{]}--{[}4{]}；对象中心追踪重建状态依赖的底层操作语言
{[}15{]}。怪异机器研究非预期计算与可利用性 {[}16{]},
{[}17{]}，不安全编译显式引入策略边界
{[}18{]}，携带证明代码工作识别证明模型之外的影子执行 {[}19{]}。有效 RPM
元数据也可产生计算，漏洞流类型可导出抽象怪异机器 {[}20{]},
{[}21{]}。不同于只展示丰富行为的工作，本文判据追问的是：在一个已接受制品的前沿上，已声明的实际计算报告是否真的因子分解未来观测商。

抽象解释提供具体/抽象、健全性和完备性词汇 {[}5{]},
{[}22{]}；本文只讨论一个前沿上唯一分配的实际计算单元。MOAT 隔离接受后的
BPF 效果 {[}23{]}，mismorphism 比较不同解释 {[}24{]}。eBPF
范围分析验证与状态嵌入针对验证器逻辑的健全性或覆盖错误 {[}25{]},
{[}26{]}。VEP 所谓``programmability''是减少验证工具链限制；Rex
通过语言级安全与语言外运行时替代独立静态验证器，处理拒绝侧的语言---验证器失配
{[}27{]}, {[}28{]}。本文 P
不同：它要求一个经原生验证器接受的固定制品解释有界族。DRACO
在验证器接受后检查功能规范；bpfverify 把 eBPF 字节码翻译为位与内存精确的
Horn 子句以进行功能验证 {[}29{]}, {[}30{]}。

2026 年预印本进一步界定边界：Heimdall
处理编译和接受后仍存的高层缺陷；Yaksha-Prashna 提取第三方 eBPF
字节码行为；bpfix
定位被拒程序中的证明丢失；信任边界语义缝隙工作分析正确语法接受后仍不足的断言
{[}31{]}--{[}34{]}。它们都未联合区分本文论断图谱的全部义务。反过来，本文
\texttt{wm\_circuit} 解释器载体也未建立 R 或 W；显式有限的冻结
stock-Linux 载体与辅助元组只在各自载体上建立 R，都不提供 P 或 W。四项
2026 工作均按预印本引用，而非声称已经同行评审。

\hypertarget{7-ux5c40ux9650ux542fux793aux4e0eux5c55ux671b}{%
\section{7.
局限、启示与展望}\label{7-ux5c40ux9650ux542fux793aux4e0eux5c55ux671b}}

\hypertarget{71-ux4e25ux683cux533aux5206ux4e0eux5206ux7c7bux8c13ux8bcd}{%
\subsection{7.1
严格区分与分类谓词}\label{71-ux4e25ux683cux533aux5206ux4e0eux5206ux7c7bux8c13ux8bcd}}

令模型

\[M=(V,I,\mathsf{Report},K_{\mathrm{obs}},P_*,\ell_*,D_R,F_R,
\mathcal A,\mathcal E,\mathsf{Drive},\mathsf{Pol})\]

还带有具名的预期语义、安全、报告抽象与粒度契约。A、C、P、R、W
分别按前文定义；P 还要求同制品门基满足 E1--E3、\(P_*\) 履行 E4-D，且
\(g\) 与常量 \(\{0,1\}\) 构成功能完备的布尔基；R
的带标签见证必须使用同一个可接纳元组。

\(\mathsf{Link}(M)\) 要求用于分类的 R 见证属于建立 P
的门/解释器执行族，并且 W
效果来自行为主体驱动同一编码计算；\(\mathsf{Unint}(M)\)
表示该解释或效果在边界预期语义之外。定义

\[\mathsf{WM}_{\mathrm{policy}}(M)=P(M)\land W(M)\land
\mathsf{Link}(M)\land\mathsf{Unint}(M).\]

\(\mathsf{Doc}(M)\)
表示同一链接行为使用已文档化语义且维持声明的安全契约；\(\mathsf{Conf}(M)\)
表示报告实现符合指定抽象；\(\mathsf{Gran}(M)\) 表示 R
碰撞由声明的报告粒度强制。更窄分类为

\[\mathsf{WM}_{\mathrm{shape}}(M)=
\mathsf{WM}_{\mathrm{policy}}(M)\land R(M)\land
\mathsf{Doc}(M)\land\mathsf{Conf}(M)\land\mathsf{Gran}(M).\]

定义直接给出 \(R\Rightarrow C\Rightarrow A\) 与
\(P\Rightarrow C\Rightarrow A\)。有限反模型表明

\[A\not\Rightarrow C,
C\not\Rightarrow P,
C\not\Rightarrow R,
P\not\Rightarrow R,
R\not\Rightarrow P,
P\not\Rightarrow W,
R\not\Rightarrow W.\]

因此裸 C、抽象解释不完备或接受不完备都不推出 P
或怪异机器分类；策略级怪异机器也可能在报告精确时成立，故一般不要求 R。

\hypertarget{72-ux9632ux5fa1ux542fux793aux4e0eux672aux6765ux5f62ux72b6ux5b9aux7406}{%
\subsection{7.2
防御启示与未来形状定理}\label{72-ux9632ux5fa1ux542fux793aux4e0eux672aux6765ux5f62ux72b6ux5b9aux7406}}

审计时应分别声明识别属性和报告，枚举行为主体可驱动的操作与环境假设，再测试报告因子分解。契约违反需要修实现；已文档化碰撞只在报告确实意图认证该关系时，才要求细化、限制或隔离。

\texttt{wm\_circuit} 载体在源码/对象和串行化前提下支持 P，但没有建立 R
或 W。显式有限的 \(M_K\) 元组针对事后声明的一步 \texttt{rac\_single}
关系和操作剪枝报告满足本文形式 R
谓词；其证据标签是轨迹局部证书，而不是一般 stock-Linux
报告结果。辅助模型 \(M_{\mathit{linux\_r\_aux\_v1}}\)
也只针对固定自定义报告和受限服务语义建立
R。由于制品与报告载体不同，\(\mathsf{Link}\) 禁止把任一 R 与解释器 P
合并。故没有已评估模型建立 \(\mathsf{WM}_{\mathrm{policy}}\) 或
\(\mathsf{WM}_{\mathrm{shape}}\)。未来形状定理必须在同一载体上推出接受的宏闭包与碰撞的报告嵌入；上下文敏感性可能破坏任一条件。

\hypertarget{8-ux7ed3ux8bba}{%
\section{8. 结论}\label{8-ux7ed3ux8bba}}

在分层的 LangSec
边界上，制品识别、接受后解释、报告因子分解、有界可编程性和策略级怪异机器状态是不同的安全判断：\(R\Rightarrow C\Rightarrow A\)
且 \(P\Rightarrow C\Rightarrow A\)；命题 3
用有限反模型否定其余列出的蕴含。\texttt{wm\_circuit} 载体记录
A，条件性见证 C，并在额外前提下支持固定 64 输入/512 门 NAND 解释器的
P，但不建立 R 或 W。辅助元组独立针对自定义报告建立 R。另一个 stock-Linux
捕获提供真实 exact-level-0
剪枝事件以及共同后缀下的不同结果；把它事后嵌入声明的有限
\texttt{rac\_single}
载体后，可得到该元组的轨迹局部、相对于操作剪枝报告的 R
证书，但不建立一般或 Linux 指定的功能报告失效。两个 R 载体都不提供 P 或
W，也不与解释器链接，故证据不建立
\(\mathsf{WM}_{\mathrm{shape}}\)、策略级怪异机器、验证器不健全、漏洞或普遍必要性定理。

\hypertarget{ux4f26ux7406ux4e0eux6570ux636eux53efux7528ux6027}{%
\section{伦理与数据可用性}\label{ux4f26ux7406ux4e0eux6570ux636eux53efux7528ux6027}}

实验在隔离本地 VM
中运行，不涉及第三方目标或生产数据路径。解释器运行没有实时内核钩子；stock-Linux
R 捕获只在隔离程序加载期间短暂把 fexit
观察器附着到验证器内部函数，观察验证器决定而不改变已接受程序的执行。任何实验都没有尝试破坏、验证器绕过或提权。

仓库为 \url{https://github.com/Emtanling/eBPF-machine}；不可变 eBPF
解释器证据位于提交
\href{https://github.com/Emtanling/eBPF-machine/tree/4309069a1f94d642d5c1402eb710e089c85059b1}{\texttt{4309069a}}
下的
\texttt{results/interpreter/interpreter-final-20260711-02/}，辅助报告实例证据位于提交
\href{https://github.com/Emtanling/eBPF-machine/tree/f665b1a2f9a772ee9b2c08a73d116ea283aa5efb}{\texttt{f665b1a}}
下的 \texttt{results/linux\_r/linux-r-v1/}。stock-Linux 轨迹证据发布在带标签的
\href{https://github.com/Emtanling/eBPF-machine/tree/V1.0/residuality-auditor/stock-linux-r-proof}{\texttt{V1.0}}
快照中，路径为
\path{residuality-auditor/stock-linux-r-proof/}。其 \texttt{MANIFEST.json}、\texttt{CHECKSUMS.sha256}、嵌入输入哈希、原始捕获、规范化证书、证明输出和检查器源码共同纳入该公开快照。随附的标准库校验器输出
\texttt{FROZEN\_PROOF\_BUNDLE\_VERIFIED}。发布前归档
\path{residuality-auditor-v0.3.0-full.tar.gz} 的 SHA-256 为
\path{5fd0a2812c8c8db2fe5508440934817c1cf9293ba0c5df31317e8b38d94a90ec}；V1.0 直接发布冻结证明载荷，不重复存储该归档。任何冻结字节变化都需要新版本、更新校验和并重新验证。

\hypertarget{ux81f4ux8c22ux4e0e-ai-ux4f7fux7528ux58f0ux660e}{%
\section{致谢与 AI
使用声明}\label{ux81f4ux8c22ux4e0e-ai-ux4f7fux7528ux58f0ux660e}}

OpenAI Codex 协助全文起草与语言修订，并在作者指导下对摘要、引言、第
5.6-\/-5.8
节和结论进行实质修订，同时协助制品代码/测试修订及检查。作者独立审查主张、证明、引用、更改与结果并承担全部责任。作者声明不存在利益冲突。

\hypertarget{ux53c2ux8003ux6587ux732e}{%
\section{参考文献}\label{ux53c2ux8003ux6587ux732e}}

为避免中文阅读版在最终投稿修订中形成第二套书目真源，本文件不重复排印完整书目；文内采用最终英文
bibliography 的 34 项显示序号，完整、投稿规范的书目信息与顺序以英文
\texttt{PAPER\_REPORT.tex} 的 \texttt{thebibliography}
为准。{[}6{]}--{[}9{]} 为
Nerode、Mealy、观测完备性与强保持理论；{[}13{]} 为 Linux v6.17
\texttt{bpf\_loop} API；{[}20{]}--{[}30{]} 为 RPM/漏洞流及 eBPF
相邻工作，其中 Rex 为 {[}28{]}、bpfverify 为
{[}30{]}；{[}31{]}--{[}34{]} 为明确标注为预印本的 2026 工作。

\end{document}
